% ═══════════════════════════════════════════════════════════════════════
\section{Livrable 2 : un examen avec la classe exam}
% ═══════════════════════════════════════════════════════════════════════

\begin{frame}{Pourquoi la classe exam}
  \small
  \begin{itemize}
    \item Des questions, sous-questions et choix multiples numérotés
      automatiquement.
    \item Des points par question ; le total et la grille de correction se
      calculent seuls.
    \item Les solutions écrites sous chaque question : un seul fichier pour
      la copie de l'élève et le corrigé.
    \item Des espaces de réponse réglables : vide, lignes, cadre.
    \item Des en-têtes et pieds de page avec « Page 1 de 3 ».
  \end{itemize}
  \vspace{4pt}
  \begin{Rem}{Un seul interrupteur}
    \small Avec \cmd{printanswers} dans le préambule, les solutions
    s'affichent ; en le mettant en commentaire (\texttt{\%}), elles
    disparaissent et l'espace de réponse est conservé.
  \end{Rem}
\end{frame}

\begin{frame}[fragile]{La structure d'un examen}
  \begin{columns}[T]
    \begin{column}{0.56\textwidth}
\begin{CodeTeX}[title={Le corps de 03-examen.tex (extrait)}]
\begin{questions}
\question[2] Lequel de ces triplets ... ?
\begin{choices}
  \CorrectChoice $3$, $4$ et $5$
  \choice $4$, $5$ et $6$
\end{choices}
\question Un triangle rectangle ...
\begin{parts}
  \part[2] Calcule l'hypoténuse.
  \begin{solution}
    $c = \sqrt{6^2 + 8^2} = \qty{10}{\centi\metre}$.
  \end{solution}
\end{parts}
\end{questions}
\end{CodeTeX}
    \end{column}
    \begin{column}{0.41\textwidth}
      \footnotesize
      \begin{itemize}
        \item \texttt{questions} : la liste de toutes les questions.
        \item \cmd{question[2]} : une question de 2 points.
        \item \texttt{parts}, \cmd{part[2]} : les sous-questions a), b).
        \item \texttt{choices}, \cmd{choice} : un choix multiple ;
          \cmd{CorrectChoice} marque la bonne réponse dans le corrigé.
        \item \texttt{solution} : la réponse, visible seulement avec
          \cmd{printanswers}.
      \end{itemize}
    \end{column}
  \end{columns}
\end{frame}

\begin{frame}[fragile]{Mettre la classe exam en français}
  \begin{columns}[T]
    \begin{column}{0.56\textwidth}
\begin{CodeTeX}[title={Préambule de 03-examen.tex (extrait)}]
\documentclass[11pt, addpoints]{exam}
% \printanswers   % corrigé : enlever le %
\pointpoints{point}{points}
\hqword{Question}
\hpword{Points}
\hsword{Résultat}
\htword{Total}
\renewcommand{\solutiontitle}{\noindent Solution : }
\pagestyle{headandfoot}
\firstpageheader{Mathématique, 3\textsuperscript{e} secondaire}
                {}{Examen --- Relation de Pythagore}
\firstpagefooter{}{Page \thepage\ de \numpages}{}
\end{CodeTeX}
    \end{column}
    \begin{column}{0.41\textwidth}
      \footnotesize
      \begin{itemize}
        \item \texttt{addpoints} : active le calcul du total.
        \item Sans ces réglages, la grille affiche Question, Points,
          Score : l'IA les oublie souvent.
        \item En-tête et pied de page : trois cases (gauche, centre,
          droite), entre accolades.
      \end{itemize}
    \end{column}
  \end{columns}
\end{frame}

\begin{frame}[fragile]{Grille, total et espaces de réponse}
\begin{CodeTeX}[title={Quelques commandes utiles}]
Cet examen compte \numquestions\ questions, pour un total de \numpoints\ points.
\gradetable[h][questions]          % grille de correction horizontale

\vspace{\stretch{1}}               % un espace qui partage la page restante
\begin{solution}[3cm]              % 3 cm réservés si le corrigé est caché
\fillwithlines{3cm}                % 3 cm de lignes pour écrire
\makeemptybox{3cm}                 % un cadre vide de 3 cm
\end{CodeTeX}
  \small Comme pour les sections, rien n'est numéroté ni additionné à la
  main : ajouter une question met à jour la numérotation, le total et la
  grille à la compilation suivante.
\end{frame}

\begin{frame}{Le rendu : copie de l'élève et corrigé}
  \begin{columns}[c]
    \begin{column}{0.36\textwidth}
      \centering
      \fbox{\includegraphics[height=0.7\textheight]{03-examen.pdf}}\\[2pt]
      \scriptsize sans \cmd{printanswers}
    \end{column}
    \begin{column}{0.36\textwidth}
      \centering
      \fbox{\includegraphics[height=0.7\textheight]{03-examen-corrige.pdf}}\\[2pt]
      \scriptsize avec \cmd{printanswers}
    \end{column}
    \begin{column}{0.25\textwidth}
      \small Le même fichier, une seule ligne de différence. La bonne
      réponse du choix multiple apparaît en gras dans le corrigé (réglage
      par défaut de la classe).
    \end{column}
  \end{columns}
\end{frame}

\begin{frame}{La consigne}
  \begin{Exercice}{Un examen sur le même thème que votre activité}
    \small
    \begin{itemize}
      \item Classe \texttt{exam}, avec l'option \texttt{addpoints} et la
        classe réglée en français.
      \item Au moins quatre questions, dont un choix multiple et une
        question à sous-questions ; des points pour chacune.
      \item Une grille de correction et le total calculés automatiquement.
      \item Une solution pour chaque question ; remettre deux PDF : la
        copie de l'élève et le corrigé.
      \item En-tête avec le niveau, pied de page « Page x de y ».
    \end{itemize}
  \end{Exercice}
\end{frame}

\begin{frame}{Prompt type : l'examen}
  \begin{Prompt}{Un examen avec la classe exam}
    \footnotesize Rédige en LaTeX un examen de \champ{3\textsuperscript{e}
    secondaire, ou d'un cours collégial} sur \champ{la relation de Pythagore}, avec la classe exam
    (option addpoints), compilable avec pdfLaTeX sur Overleaf.

    Contenu : \champ{5} questions pour un total de \champ{20} points ;
    au moins un choix multiple (choices, CorrectChoice), une question avec
    sous-questions (parts), et un problème en contexte réel. Mets une
    solution détaillée dans un environnement solution sous chaque question,
    et laisse un espace de réponse adapté.

    Mise en forme : la classe exam entièrement en français (pointpoints,
    hqword, hpword, hsword, htword, solutiontitle) ; une grille de
    correction horizontale ; en-tête avec le niveau et le titre ; pied de
    page « Page x de y ». Laisse \textbackslash printanswers en
    commentaire au début du fichier. Utilise siunitx avec virgule décimale
    pour les nombres et les unités. Donne un seul fichier .tex.
  \end{Prompt}
\end{frame}

\begin{frame}{Prompts types : réutiliser ses questions}
  \begin{columns}[T]
    \begin{column}{0.49\textwidth}
      \begin{Prompt}{Une version B}
        \scriptsize Voici mon examen en LaTeX : \champ{fichier}. Produis
        une version B de même difficulté : mêmes questions, mêmes points,
        mais d'autres valeurs numériques, choisies pour que les réponses
        restent simples. Recalcule toutes les solutions et montre-moi les
        calculs pour que je les vérifie.
      \end{Prompt}
    \end{column}
    \begin{column}{0.49\textwidth}
      \begin{Prompt}{Ajouter une question}
        \scriptsize Ajoute à mon examen une question de \champ{4} points
        sur \champ{la réciproque de la relation de Pythagore}, avec deux
        sous-questions et leur solution. Donne seulement le bloc
        \textbackslash question à insérer, dans le même style que les
        autres.
      \end{Prompt}
    \end{column}
  \end{columns}
  \vspace{6pt}
  \small Dans les deux cas : refaire soi-même les calculs du corrigé avant
  de distribuer la copie.
\end{frame}
